\section{The two requested Lean targets}\label{sec:formal}

\subsection{Yield conservation: a core theorem already exists}
The file
\art{crypto:research/factor_base_designs_20261015/lean/FactorBaseDesigns/Yield.lean}
defines \art{count D phi g} as the length of the filtered list of domain
entries mapping to $g$. Its theorem \art{sum_count} states, exactly in
mathematical notation,
\[
 \text{$Gs$ has no duplicates and contains every $\varphi(x)$ for $x\in D$}
 \quad\Longrightarrow\quad
 \sum_{g\in Gs}\operatorname{count}(D,\varphi,g)=\operatorname{length}(D).
\]
It requires decidable equality of the target type, not a group law.
\art{patterns_neutral} compares two equally long domains;
\art{restrict_le} proves that filtering the domain cannot increase
any target's multiplicity. The committed axiom log lists
\art{propext} and \art{Quot.sound} for these three declarations. A
committed build log reports a successful build. The present campaign also
rebuilt this nine-job project with its pinned Lean 4.34.1 toolchain and
reran its axiom audit; fresh logs are retained under
\art{crypto:research/ideas_execution_20261018/runs/}.

The campaign's new \art{IdeasExecution.OrderedYield} module enumerates the
ordered Cartesian power and proves \art{tuples_length},
\art{ordered_total_yield} and \art{ordered_patterns_neutral}. Its build
passes; its separate axiom audit lists only \art{propext} and
\art{Quot.sound}. This is an application bridge for ordered list-domain
multiplicity, not a new conservation principle.
For the multiset convention one still needs the enumeration length
$\binom{B+m-1}{m}$ without duplicate multiset entries. Prove that the
sum map and full target enumeration meet \art{sum_count}'s hypotheses.
For a mean in $\mathbb Q$, require a nonempty target set before dividing.
Those remaining bridges instantiate the existing theorem rather than create
a new conservation principle.

\subsection{Compact-pipeline genericity: an interface bridge is missing}
The exact theorem \art{GenericInvariance.run_map} in
\art{crypto:research/conductor_gap_cm_pairs_20261010/lean/GenericInvariance.lean}
states that an injective map preserving the declared multiplication and
inversion operations commutes with every bounded run of its handle-only
algorithm. \art{dlog_invariant} specializes it to a two-handle DLP input.
The source contains no elliptic-curve construction and no probabilistic
DLP lower bound. Its explicit representation-hashing limitation is
binding, as discussed above.
The campaign also rebuilt this four-job project and reran its axiom audit;
successful checking does not extend its interface to representation hashing.

A faithful target for the compact pipeline has three layers:
\begin{enumerate}
\item Specify a concrete fragment with generically constructed advice,
known scalar actions, membership through stored equality, and all index
operations exposed. Prove an abstraction relation between concrete
states and symbolic states across every executed instruction.
\item Prove a transcript simulation, including target-independent setup,
lookup/multiplicity semantics, exceptional cases and the exact resource
counters. Distinguish random-label hashing from the stricter
handle-only theorem rather than extending that theorem by prose.
\item Import or separately formalize the appropriate generic
preprocessing lower bound with its success probability, advice units,
automorphism assumptions and parameter range. It is not a corollary of
injective transport invariance alone.
\end{enumerate}
No inspected Lean file closes all three layers for the current compact
pipeline. A machine check of a small rewritten toy algorithm would not
close the concrete-code correspondence. This target therefore remains
open. The draft does not insert an unproved axiom or a \art{sorry} and
call that a formalization.
